\begin{solution}
    \begin{equation*}
    \begin{aligned}
        \textbf{Dados:}\quad 
            g(t):=t^{6}-1, \qquad
            f(t):=t^{4}-1, \qquad
            \bigl(g,f\in K[t]\bigr)
    \end{aligned}
    \end{equation*}
        
        Como el grado de $g(t)$ es mayor que el de $f(t)$, entonces por el algoritmo de Euclides sabemos que el máximo común divisor entre estos dos polinomios satisface $d(t)= \gcd(g(t),f(t)) = \gcd(f(t),r(t))$, donde $r(t)$ es el resto de la división de $g(t)$ entre $f(t)$. Por lo que realizando la division de polinomios:
      \[
      \begin{tikzpicture}[baseline={(current bounding box.north)}]
        \node at (0,0.3) {$t^6$};
        \node at (0,-0.3) {$-t^6$};
        \node at (0,-.8) {$\underline{\hspace{1.5cm}}$};
        \node at (.6,-.3) {$+$};
        \node at (.9,-.3) {$\phantom{-}t^2$};

        \node at (.5,0.3) {$-$};
        \node at (1,0.3) {$1$};
        \draw (1.4,0.6) -- (3,0.6);
        \draw (1.4,0.6) -- (1.4,-1.3);
        \node at (2,1.0) {$t^4 - 1$};
        \node at (1.8,0.3) {$t^2$};
        \node at (0,-1.2) {$t^2$};
        \node at (0.5,-1.2) {$-$};
        \node at (1,-1.2) {$1$};
      \end{tikzpicture}
      \]
        Es decir que:

        \[
        \begin{aligned}
        g(t)&= t^2 (t^4-1) + (t^2-1)\\
            &= t^{2}\,f(t)+\underbrace{\bigl(t^{2}-1\bigr)}_{=:r_{1}(t)} \qquad \deg r_{1}<\deg f\\
        \end{aligned}
        \]

        Luego $\gcd(g(t),f(t))=\gcd(f(t),r_{1}(t))$ y continuamos con la división de $f(t)$ entre $r_{1}(t)$:

        \[
          \begin{tikzpicture}[baseline={(current bounding box.north)}]
            \node at (0,0.3) {$t^4$};
            \node at (0,-0.3) {$-t^4$};
            \node at (0,-.8) {$\underline{\hspace{2.4cm}}$};
            \node at (.6,-.3) {$+$};
            \node at (.9,-.3) {$\phantom{-}t^2$};
    
            \node at (.5,0.3) {$-$};
            \node at (1,0.3) {$1$};
            \draw (1.4,0.6) -- (3,0.6);
            \draw (1.4,0.6) -- (1.4,-2.6);
            \node at (2,1.0) {$t^2- 1$};
            \node at (2.1,0.3) {$t^2+1$};
            \node at (0,-1.2) {$t^2$};
            \node at (0.5,-1.2) {$-$};
            \node at (1,-1.2) {$1$};
            \node at (0,-2.2) {$\underline{\hspace{2.4cm}}$};
            \node at (0,-1.8) {$-t^2$};
            \node at (0.6,-1.8) {$+$};
            \node at (1,-1.8) {$1$};
            \node at (1,-2.6) {$0$};

          \end{tikzpicture}
          \]
        Es decir:
        \begin{align*}
          f(t)&=\bigl(t^{2}+1\bigr)\,r_{1}(t)+\underbrace{0}_{r_{2}(t)}
        \end{align*}
        
        Por lo tanto $\gcd(g(t),f(t))=\gcd(r_{1}(t),0)=\gcd(r_{1}(t))=r_{1}(t)$. Y dado que de la primera división:
        \[
        r_{1}(t)=t^{2}-1=g(t)-t^{2}\,f(t).
        \]
        
        Por lo tanto existen
        \[
        a(t)=-t^{2}, \qquad b(t)=1 \quad\bigl(a(t),b(t)\in K[t]\bigr)
        \]
        tales que
        \[
        \boxed{\;
        \gcd\!\bigl(g(t),f(t)\bigr)=t^{2}-1
            =a(t)\,f(t)+b(t)\,g(t)
            =\bigl(-t^{2}\bigr)\,f(t)+1\cdot g(t)
        \;} .
        \]
\end{solution}